\documentclass{article}
\usepackage{hyperref}
\usepackage{color}
\usepackage{enumitem}

 \oddsidemargin 0pt % -10.4mm
 \evensidemargin 0pt % -10.4mm
 % \marginparwidth 0pts
 \marginparsep 0pt
 \topmargin 0mm % -10.4mm
 \headheight 0pt
 \headsep 0pt
 \textheight 232mm % 257.2mm
 \textwidth 165mm % 169.2mm
 \footskip 25pt %


\begin{document}

\begin{center}
\rule{\textwidth}{.0075in} \\
\rule[3mm]{\textwidth}{.0075in}\\

Universidade do Algarve\hfill Introdução à Cybersegurança\hfill 2026\\[3ex]

{\Large\bf HW4} \\[3ex]

\vspace*{0.30cm} Nome: 
  \raisebox{.2ex}{\makebox[0pt][l]{}}{\bf{}Nuno Carmo Vivas} \ %
  \textnumero:  \;\textbf{74799}

\vspace*{0.30cm} Nome: 
  \raisebox{.2ex}{\makebox[0pt][l]{}}{\bf{}Francisco Santos} \ %
  \textnumero:  \;\textbf{90118}

\rule{\textwidth}{.0075in} \\
\rule[3mm]{\textwidth}{.0075in} \\
\end{center}

\bigskip

\begin{center}
\setlength{\fboxrule}{2pt}
  \fbox{\parbox{0.8\linewidth}{
      \begin{itemize}
      \item Somente um aluno por grupo deve submeter o relatório.
      \item Utilize Overleaf para editar e produzir o PDF do relátorio.
      \end{itemize}  
  }}
\end{center}

\medskip

\section*{Introducção}
Devem trabalhar em grupos de 2 alunos para este laboratório. O
objetivo deste laboratório é, em primeiro lugar, de comprenderem como são registadas
as vulnerabilidades de software e, em segundo lugar, de escolherem uma vulnerabilidade
de software, ou CVE, para analisar. A vulnerabilidade de software deve ser de alta
prioridade como score maior do que 9,5 e ter um exploit disponível. Um exploit é um programa que fornece
instruções passo a passo sobre como tornar o software em questão vulnerável a um
ciberataque. Neste laboratório, deverão escrever e enviar um relatório sobre a CVE
escolhida.

O primeiro passo é familiarizar-se com o seguinte website:
\href{https://www.cvedetails.com/} {https://\-www.\-cve\-de\-tails.\-com/}, que mantém um
registo de todas as CVEs (vulnerabilidades) de software registadas na base de
dados. À esquerda, encontrará uma opção de ``Search'' (pesquisa) para procurar CVEs
de acordo com os seus próprios critérios de pesquisa. Na secção CVSS Score, deve
escolher um mínimo de 9.5 e um máximo de 10.0 (prioridade crítica) e selecionar a
opção que diz ``Has public or potential exploit''. Pode escolher os outros
parâmetros de acordo com os seus interesses específicos, como o produto, datas,
categorias de vulnerabilidade, etc. O objetivo desta etapa é selecionar uma CVE
específica sobre a qual irá escrever o seu relatório.

A segunda etapa é familiarizar-se com as CWEs. Trata-se de categorias mais amplas
ou gerais de vulnerabilidades de software: uma CWE engloba ou inclui várias
CVEs. Visite este site: \href{https://cwe.mitre.org/} {https://cwe.mitre.org/}. A
CVE que escolheu na primeira etapa está ligada a uma CWE. Explore esta CWE nesta
segunda etapa.

\section*{Conteúdo do Relatório}

Deve enviar um relatório em formato PDF que inclua as seguintes secções: 

\medskip

\noindent
\textbf{Secção 1 (30\%)}: Análise da falha (``flaw analysis''). Explique a natureza da
vulnerabilidade escolhida e quais os princípios de design de software seguro e/ou de
codificação que foram violados.

\begin{center}
\setlength{\fboxrule}{1.5pt}
\fbox{\parbox{0.8\linewidth}{
    \begin{center}
      \begin{tabular}{c}
        {\bf Dicas}
      \end{tabular}
    \end{center}  
      \begin{itemize}
      \item Discutimos os princípios de design de software na primeira
        e segunda semanas; consulte os slides na Tutoria. 
      \item Deve mapear a sua CVE num ou mais princípios de
        design de software e explicar por que razão a falha discutida
        na sua CVE viola esses princípios.
      \item Deve também mapear as falhas encontradas pelo seu CVE na tríade CIA.   
     \end{itemize}  
  }}
\end{center}


\medskip

\noindent
\textbf{Secção 2 (30\%)}: Exploits e Violações de Código (``code
exploits and violations''). Descreva e demonstre o código exploit que
leva a uma violação de segurança no produto/componente de
software. Mostre e descreva os passos que o exploit utiliza para
cometer a violação. Apresente exemplos das violações de software.

\begin{center}
\setlength{\fboxrule}{1.5pt}
\fbox{\parbox{0.8\linewidth}{
    \begin{center}
      \begin{tabular}{c}
        {\bf Dicas}
      \end{tabular}
    \end{center}  
      \begin{itemize}
      \item Certifique-se de que consegue executar o exploit e reproduzi-lo. 
      \item Por definição, um exploit inclui passos e refere-se a um
        código que expõe uma vulnerabilidade.
      \item Se pretende executar o exploit localmente no seu
        computador, poderá querer pesquisar na internet por "chroot",
        "chroot jail" ou termos semelhantes.   
     \end{itemize}  
  }}
\end{center}


\noindent
\textbf{Secção 3 (30\%)}: Remediação (``Remediation''). Como pode a vulnerabilidade ser remediada? Explique
como poderia ser corrigido o problema, por exemplo, utilizando uma correção de
código. Quê outras alternativas seriam viáveis?


\begin{center}
\setlength{\fboxrule}{1.5pt}
\fbox{\parbox{0.8\linewidth}{
    \begin{center}
      \begin{tabular}{c}
        {\bf Dicas}
      \end{tabular}
    \end{center}  
      \begin{itemize}
      \item As remediações podem assumir diferentes formas, como
        ``patches'' (remendos), atualizações, etc. Certifique-se de que funcionam.   
     \end{itemize}  
  }}
\end{center} 


\noindent
\textbf{Secção 4 (5\%)}: Vulnerabilidades Relacionadas. Relacionar a vulnerabilidade
apresentada com outras vulnerabilidades existentes, descrevendo as suas semelhanças
e diferenças, e como se relacionam. Liste o CVE-ID de cada vulnerabilidade
relacionada.

\begin{center}
\setlength{\fboxrule}{1.5pt}
\fbox{\parbox{0.8\linewidth}{
    \begin{center}
      \begin{tabular}{c}
        {\bf Dicas}
      \end{tabular}
    \end{center}  
      \begin{itemize}
      \item Explique porquê e como é que as CVEs se relacionam com a sua própria CVE.   
     \end{itemize}  
  }}
\end{center} 

\noindent
\textbf{Secção 5 (5\%)}: Referências. De onde tomou (pegou) a informação relatada neste
relatório? 



\section*{Secção 1: Análise da Falha}



\section*{Secção 2: Exploits e Violações de Código}



\section*{Secção 3: Remediação}



\section*{Secção 4: Vulnerabilidades Relacionadas}



\section*{Secção 5: Referências}




\end{document}
